\section{Smooth tests for the regularized commutator under a fixed zero gap}
\label{sec:regularized-smooth}

\paragraph{MA-RS-004 (arbitrary compact smooth tests under a fixed horizontal gap).}
Retain the actual closed multiplier operators of MA-RS-001--003.
Suppose $1/2<\sigma<1$ and, for some $\delta>0$,
\[
       |\sigma-\operatorname{Re}\rho|\ge\delta
          \quad\text{for every nontrivial zero of }\zeta.
 \tag{MA-RS.32}
\]
Put $a=\sigma-1/2$. Then the conclusions of MA-RS-002--003
hold for every $g\in C_c^\infty(\mathbb R)$ and every bounded
smooth $\phi$ with
\[
       \phi(-\infty)=1,\qquad \phi(+\infty)=0,
                         \qquad\phi'\in C_c^\infty(\mathbb R).
 \tag{MA-RS.33}
\]
In particular the actual composition on the compact-frequency core
\[
             C_gZ_{a,S}[J_{a,S},M_\phi]
 \tag{MA-RS.34}
\]
extends uniquely to a trace-class operator on
$\mathcal H_a=L^2(e^{2ax}dx)$, with the exact kernel and trace
in (MA-RS.16)--(MA-RS.17). Its residue formula is still
(MA-RS.26), and its reflected pair still has (MA-RS.28).
The map $g\mapsto K_{a,S,\phi}(g)$ is continuous into trace
class on each fixed compact test-support space.

Under the stipulated Dirichlet input $A$, applied to rational
zeta, (MA-RS.32) holds for every $1/2<\sigma<1$.
For these lines the zero correction is absent, and the formulas
become
\[
 \begin{aligned}
 \operatorname{Tr}K_{a,S,\phi}(g)
     &=P_{\overline S}^{+}(g)-\mathcal Mg(1/2),\\
 \operatorname{Tr}\mathbb K_{a,S,\phi}(g)
     &=B_{\overline S}(g)
          -\mathcal Mg(1/2)-\mathcal Mg(-1/2).
 \end{aligned}
 \tag{MA-RS.35}
\]
Thus a smoothed commutator can be well defined and trace class
even though the bare M\"obius inverse has no bounded extension
between matching narrow weighted norms, as MA-MO-001 proves.

\begin{proof}
We first improve the inverse-multiplier estimate, keeping its
zero-location premise explicit. Use the finite-disk proof in
MA-RS-001, centered at $2+it$ with a zero-free boundary radius
$12<R<13$ for large $|t|$. Its factorization $\zeta=G B$
has $M=O(\log(2+|t|))$ zeros in the disk. Harnack's inequality
and the fixed-strip polynomial upper bound give
$|G(w)|\ge c(2+|t|)^{-C}$ on the inner half disk.

At $w=\sigma+it$ each finite Blaschke factor has modulus at
least $\delta/(2R)$, by (MA-RS.32). Replace this lower bound
by its minimum with $1/2$ if necessary. Its product over the
$M$ factors is at least a fixed negative power of $2+|t|$.
The same estimate is uniform on a fixed disk of radius
$\epsilon>0$ about $\sigma+it$, where $\epsilon$ is smaller
than $\delta/2$ and than the distance of the line from the
pole at $1$. This disk contains no zero or pole.
Cauchy's estimate for the reciprocal on a slightly smaller
disk consequently proves, for every integer $r\ge0$,
\[
 \left|\frac{d^r}{dt^r}\frac1{\zeta_S(\sigma+it)}\right|
       \le C_r(2+|t|)^{D_r}.
 \tag{MA-RS.36}
\]
The finitely many removed Euler factors and their inverses
are bounded with all derivatives on this fixed narrow strip.
The compact range of $t$ follows by continuity and the absence
of zeros on the line. All derivatives of
$A(t)=\zeta_S(\sigma+it)$ have polynomial growth as before.
This uses the actual gap (MA-RS.32), not merely a zero-free
line without quantitative control.

For compact smooth $g$ and $\phi'$, the functions
\[
     m(t)=\mathcal Mg(a+it),\qquad
     p(t)=\int e^{itx}\phi'(x)\,dx
\]
belong to $\mathcal S(\mathbb R)$, with continuous Schwartz
seminorms on each fixed compact support space. In the unitary
weighted Fourier picture, multiplication by $\phi$ preserves
the Schwartz space: $\phi$ is bounded and every positive-order
derivative is compactly supported and smooth. The multipliers
$A$ and $j=A^{-1}$ preserve Schwartz functions by their
polynomial derivative bounds. Thus, for every compact-frequency
core vector, both $J_{a,S}M_\phi$ and
$Z_{a,S}M_\phi J_{a,S}$ are defined, and their required
images lie in the domain of the next operator. This checks
the actual composition (MA-RS.34), rather than defining
it by a trace.

The exact distributional commutator calculation of MA-RS-002
then gives the kernel
\[
       \frac{i}{2\pi}m(t)A(t)
                   \frac{j(t)-j(u)}{t-u}\,p(t-u).
 \tag{MA-RS.37}
\]
Every mixed derivative of the divided difference has polynomial
growth in $(t,u)$: differentiate
$\int_0^1j'(u+\theta(t-u))\,d\theta$ and use
(MA-RS.36). The linear change of variables $(t,u)\mapsto
(t,t-u)$ shows that the product of any such polynomial with
any derivatives of $m(t)p(t-u)$ decreases faster than every
power in both variables. Hence (MA-RS.37) is a Schwartz
kernel. The Hermite-expansion proof in MA-RS-002 gives trace
class, the diagonal-integral formula, and continuity in trace
norm from finitely many Schwartz seminorms.

The contour argument of MA-RS-003 also applies unchanged.
At its selected horizontal heights the logarithmic derivative
has polynomial growth, while $\mathcal Mg$ decreases faster
than every power on the fixed strip. The horizontal integrals
therefore tend to zero, and the same residues give
(MA-RS.26) and (MA-RS.28). No Gevrey estimate is needed once
(MA-RS.36) is available.

Under $A$ all the nontrivial zeta zeros have real part $1/2$.
For each fixed $\sigma>1/2$ use $\delta=\sigma-1/2$.
There are no crossed nontrivial zeros on the right of the
line, nor outside its reflected pair. The pole at $1$
is still crossed and gives exactly (MA-RS.35).
\end{proof}

\paragraph{A concrete cutoff distinction, still present under $A$.}
Choose $g\ge0$, $g\ne0$, with compact support in
$(0,\log2)$. Both prime sums in (MA-RS.35) vanish,
regardless of the finite set $S$. Nevertheless
\[
 \operatorname{Tr}K_{a,S,\phi}(g)=-\mathcal Mg(1/2)<0,
 \qquad
 \operatorname{Tr}\mathbb K_{a,S,\phi}(g)
        =-\mathcal Mg(1/2)-\mathcal Mg(-1/2)<0.
 \tag{MA-RS.38}
\]
If the cutoff itself is compactly supported, the same
Schwartz-kernel proof applies but $p(0)=\int\phi'=0$,
so its commutator trace is zero. Such compact-cutoff
commutators cannot converge in trace norm to either
nonzero-trace operator in (MA-RS.38). This is an actual
operator distinction and a necessary pole correction;
it is not a refutation of a formula that already includes
the two polar terms.

\paragraph{Scope of this improvement.}
The unconditional generic-line result MA-RS-001--003 retains
its stated Gevrey scope. MA-RS-004 uses the additional fixed
horizontal gap and covers all compact smooth tests. For the
rational zeta sectors that gap follows from $A$; this use of
zero location is explicit. The construction concerns the actual
one-sided closed multipliers and their smoothed commutator.
It does not by itself identify them with the orthogonal
Paley--Wiener projections of Masullo's global weighted space.
